\documentclass[a4paper,12pt]{amsart}
\usepackage[french]{babel}
\usepackage[utf8x]{inputenc}
\usepackage{graphicx}
\usepackage{enumitem}
\usepackage{amssymb}
\usepackage{geometry}

\newcommand{\C}{\ensuremath{\mathbb{C}}}
\title{Fonctions de la variable complexe - Correction}
\begin{document}
\maketitle

\section*{Exercice}
\begin{enumerate}
\item Soit $\mathcal{C}_r \colon t \in [0, 2\pi] \mapsto z_0 + r
  \exp(it)$. Soit $V$ un voisinage de $z_0$ sur lequel $f$ est bornée
  en module par un réel $M > 0$. Il existe $r_0 > 0$ tel que pour tout
  $r \leq r_0$, le cercle $\mathcal{C}_r$ soit dans $V$. On a alors,
  pour tout entier $n \geq 1$ et tout $u \in \mathcal{C}_r$:
\[
|(u-z_0)(u-z_0)^{n-1}f(u)| \leq r^n M
\]
Le lemme de Jordan donne:
\[
\lim_{r \to 0} \int_{\mathcal{C}_r} (u-z_0)^{n-1} f(u) du = 0
\]
Comme par ailleurs:
\[
 \int_{\mathcal{C}_r} (u-z_0)^{n-1} f(u) du = i 2 \pi b_n
\]
on a bien $b_n=0, n \geq 1$.
\item L'application $h$ est bornée par $\epsilon^{-1}$. Elle admet
  donc un prolongement analytique en $z_0$ en vertu de la question
  précédente.
\item $h$ ne peut s'annuler qu'en $z_0$. Si elle ne s'annule pas, on a
  $f(z_0) = h(z_0)^{-1}+a$ et $z_0$ n'est pas une singularité de
  $f$. Sinon, $h$ admet au voisinage de $z_0$ un développement de la forme:
\[
\sum_{n \geq p}c_n(z-z_0)^n
\]
avec $c_p \neq 0$. On en déduit que $(z-z_0)^ph^{-1}(z)$ est
holomorphe au voisinage de $z_0$ et donc que $z_0$ est pôle d'ordre
$p$ de $f$.
\end{enumerate}
\end{document}